\section{Operator Economics}
\label{sec:operator-economics}

\Zoo{} D-Chain validators stake \texttt{\$ZOO}; rewards are paid as
a share of trading-fee revenue plus \texttt{\$AI} from compute proofs
delegated to the \Hanzo{} AI Chain. This section documents the
mechanics.

\subsection{Validator stake}

\Zoo{} D-Chain validators bond \texttt{\$ZOO} as their primary
stake. The minimum stake is set by the \Zoo{} \DAO{} and revisited
on a quarterly governance cadence; at activation it is denominated
in a \texttt{\$ZOO} amount sufficient to align validator incentives
with chain integrity but low enough that small-scale validators can
participate.

Optional secondary stake of \texttt{\$AI} is permitted; \texttt{\$AI}
stake gives the validator priority access to AI mining-receipt
attribution from the \Hanzo{} AI Chain inference flow that anchors
through A-Chain.

\subsection{Reward sources}

Validators earn from three sources:

\begin{enumerate}[leftmargin=1.4em]
  \item \textbf{Trading-fee share (\texttt{\$ZOO} or fee-token
        denominated).} A configurable share of taker fees on every
        trade flows to validators. Default split at activation:
        $40\%$ validator pool, $30\%$ protocol treasury, $20\%$ maker
        rebate, $10\%$ conservation allocation.
  \item \textbf{Block production rewards (\texttt{\$ZOO}).}
        Per-block subsidy that decays on a defined schedule. The
        subsidy is denominated in \texttt{\$ZOO} drawn from the
        treasury per the \Zoo{} \DAO{} schedule.
  \item \textbf{AI mining receipts (\texttt{\$AI}).}
        Validators that operate or delegate to \Hanzo{} AI Chain
        compute capacity receive a portion of inference receipts as
        \texttt{\$AI}. The mining loop binds compute to receipts
        through TEE attestation (see \cite{zoo-4-0} \S AI mining).
\end{enumerate}

\subsection{Slashing}

Slashing is enforced via Q-Chain ceremonies under LP-105
\cite{lp-105} for double-sign and via A-Chain attestation gaps for
liveness failures. Specific slashable events:

\begin{itemize}[leftmargin=1.2em]
  \item \textbf{Double-sign.} Signing two conflicting blocks at the
        same height. Severe penalty: significant percentage of
        bonded \texttt{\$ZOO} burned plus removal from validator
        set.
  \item \textbf{Liveness failure.} Missing more than the threshold
        number of blocks within an epoch. Moderate penalty:
        per-block missed-attestation deduction.
  \item \textbf{Compromised key not rotated.} Failing to rotate a
        compromised key within the LP-017 schedule
        \cite{lp-134} after a published advisory. Penalty: removal
        from validator set.
  \item \textbf{Compliance precompile bypass attempt.} Submitting
        blocks containing trades that bypass the \Liquidity{}
        precompile attestation. Severe penalty: full bonded stake
        burned plus permanent ejection.
\end{itemize}

The compliance-precompile-bypass slash is the most severe because
the soundness of the \Liquidity{} Protocol integration
\cite{liquidity-protocol-formal} depends on the precompile being
the only mutator of \texttt{SecurityClass} state. Validators that
attempt to bypass it break the soundness contract, so the slash is
total.

\subsection{Maker incentives}

A portion of taker fees flows back to makers as a per-fill rebate.
The default maker rebate at activation is $1$ bp on $5$ bp taker
fee (so the maker receives $1$ bp per share traded against their
resting order; the venue keeps $4$ bps for the validator pool,
treasury, and conservation allocation). Per-pair rebates can be
configured by \Zoo{} \DAO{} governance to attract liquidity in
specific securities.

\subsection{Hanzo AI Chain compute integration}

A validator running \Zen{}-family inference workloads inside a
TEE-attested container can earn \texttt{\$AI} simultaneously with
running \Zoo{} D-Chain block production. The hardware overlap is
deliberate: the \Quasar{}\textsc{gpu} substrate \cite{lp-132}
expects validators to run on data-center \GPU{}s with confidential
compute; those same \GPU{}s execute \Zen{} inference. A single
validator hardware investment serves two revenue streams.

The economic relationship to the \Hanzo{} AI Chain is documented
in the \Hanzo{} AI Chain whitepaper; \Zoo{} validators participate
through delegation and through co-location of inference workers
on the validator hardware.

\subsection{Protocol treasury}

The protocol treasury is funded from $30\%$ of taker fees by default.
It is governed by the Zoo \DAO{} and supports:

\begin{itemize}[leftmargin=1.2em]
  \item Ongoing protocol development (engineering, audits, security
        reviews).
  \item Listing-onboarding subsidies for new asset classes.
  \item Liquidity bootstrapping for newly-listed pairs.
  \item Insurance fund against validator slashing edge cases that
        could degrade end-user experience.
\end{itemize}

\subsection{Conservation allocation}

Consistent with \Zoo{} Labs Foundation's 501(c)(3) mission, $10\%$
of taker fees flow to a conservation allocation. The allocation is
administered by the Zoo Foundation under the same impact-metric
oracle and reporting framework as the rest of the foundation's
on-chain conservation grants. This is a structural property of the
fee schedule, not a discretionary policy: every trade contributes
directly to species preservation, with on-chain reporting available
for audit.

\subsection{Validator hardware budget}

\begin{itemize}[leftmargin=1.2em]
  \item Data-center \GPU{} (H100-class or successor) with
        confidential compute support.
  \item NVMe storage for state plus order-book persistence.
  \item Network: high-bandwidth, low-latency (sub-millisecond to
        peer set; sub-100ms to global edge).
  \item Operational redundancy: hot standby, key-rotation tooling,
        TEE-attestation tooling.
\end{itemize}

The hardware budget is consistent with \Lux{} validator
expectations under \Quasar{} 4.0 (LP-105 \cite{lp-105}), which
is intentional: the \Zoo{} D-Chain validator does not have to be
a different hardware class from the \Lux{} mainnet validator.

\subsection{Economic security at activation}

At activation (April 20, 2026), the validator set is sized to
provide sufficient stake-weighted security against a single-actor
takeover; the threshold for cert finality is $t \geq \lceil 2n/3
\rceil$. The activation validator-set size and per-validator
minimum stake are set by the \Zoo{} \DAO{} prior to activation;
the parameters are public on-chain.
